
Der vorherige Abschnitt hat gezeigt, dass ERP-Systeme KMU dabei unterstützen
können, knappe Ressourcen wirksamer einzusetzen. 

Gleichzeitig entsteht daraus ein zentrales Spannungsverhältnis. ERP-Systeme
sollen Daten, Prozesse und Verantwortlichkeiten strukturieren, setzen aber für
ihre Einführung bereits ein Mindestmaß an Datenqualität, Prozessklarheit und
organisatorischer Vorbereitung voraus \parencite{Jha2008,Leyh2015}. Dieses Spannungsverhältnis kann als
ERP-Paradoxon verstanden werden: Das System soll Fähigkeiten stärken, die für
seine erfolgreiche Einführung bereits teilweise vorhanden sein müssen.

Dieses Paradoxon betrifft grundsätzlich auch Großunternehmen. Auch dort
erfordert eine ERP-Einführung belastbare Daten, nachvollziehbare Prozesse,
klare Zuständigkeiten sowie Zeit und finanzielle Mittel \parencite{Leyh2015}. Im Kontext von KMU
tritt dieses Spannungsverhältnis jedoch besonders deutlich hervor, da gerade
Ressourcen, Prozessklarheit und Verantwortlichkeiten häufig weniger stark
ausgeprägt sind als in Großunternehmen \parencite{Leyh2015,Jha2008,Nach2008}. Dadurch können
die Voraussetzungen, die für eine erfolgreiche ERP-Einführung benötigt werden,
bei KMU schneller zu einer eigenen Hürde werden.

Dies zeigt sich zunächst bei den Daten. ERP-Systeme sollen eine gemeinsame
Datenbasis schaffen und Informationen unternehmensweit verfügbar machen. Für
die Einführung ist jedoch bereits eine gewisse Qualität der vorhandenen Daten
notwendig \parencite{Leyh2017,Yusuf2024}. Fehlerhafte, unvollständige oder uneinheitliche Daten erschweren die
Übernahme in ein integriertes System und können die spätere Nutzung
beeinträchtigen \parencite{Jha2008}. 

Ähnlich verhält es sich mit den Geschäftsprozessen. ERP-Systeme sollen Abläufe
standardisieren, transparenter machen und bereichsübergreifend verbinden. Damit
dies gelingen kann, müssen bestehende Prozesse jedoch bereits in ausreichendem
Maß bekannt und nachvollziehbar sein \parencite{Leyh2015,Alaskari2021}. Nur wenn klar ist, wie zentrale Abläufe
funktionieren, können sie sinnvoll in einem ERP-System abgebildet oder
angepasst werden \parencite{Leyh2015}.

Auch organisatorische Verantwortlichkeiten müssen zumindest grundlegend
geklärt sein. ERP-Projekte betreffen unterschiedliche Unternehmensbereiche und
erfordern Entscheidungen über Prozesse, Daten und Systemnutzung. Dafür braucht
es klare Zuständigkeiten und Entscheidungswege \parencite{Leyh2017,Leyh2015}. Zudem
benötigt die Einführung Zeit und finanzielle Mittel, obwohl ERP-Systeme gerade
unter Ressourcenknappheit eingeführt werden, um Ressourcen langfristig
effizienter zu nutzen \parencite{Leyh2017}.

Das ERP-Paradoxon besteht somit darin, dass ERP-Systeme in KMU zwar helfen
können, Ordnung, Transparenz und Effizienz aufzubauen, dafür aber bereits ein
Mindestmaß an Ordnung, Transparenz und Ressourcen voraussetzen. Während diese
Anforderung grundsätzlich für Unternehmen jeder Größe gilt, ist sie für KMU
besonders kritisch, weil begrenzte Ressourcen, weniger formalisierte Prozesse
und unklarere Verantwortlichkeiten die Erfüllung dieser Voraussetzungen
zusätzlich erschweren können.
